\chapter[Tutorato 1 --- 8 ottobre 2026]{Tutorato 1: fattoriale e dimostrazioni per induzione}
\label[appendix]{chap:tutorato-01}

\begin{exercise}[Semplificazione di un rapporto di fattoriali]
\label{ex:tutorato01-rapporto-fattoriali}
Per \(n\in\N\), semplificare l'espressione
\[
  \frac{(n+2)!+(n+1)!}{(n+3)!}.
\]
\end{exercise}

\begin{proof}[Soluzione dell'esercizio]
Per la \cref{def:fattoriale},
\[
  (n+2)!=(n+2)(n+1)!,
  \qquad
  (n+3)!=(n+3)(n+2)(n+1)!.
\]
Raccogliendo \((n+1)!\) al numeratore, si ottiene
\begin{align*}
  \frac{(n+2)!+(n+1)!}{(n+3)!}
    &=\frac{\bigl((n+2)+1\bigr)(n+1)!}
      {(n+3)(n+2)(n+1)!}\\
    &=\frac{(n+3)(n+1)!}{(n+3)(n+2)(n+1)!}\\
    &=\frac{1}{n+2}.
\end{align*}
Il denominatore iniziale e i fattori semplificati sono strettamente
positivi per ogni \(n\in\N\); tutte le divisioni sono quindi lecite,
anche nel caso \(n=0\).
\end{proof}

\begin{exercise}[Somma dei primi \(n\) numeri dispari positivi]
\label{ex:tutorato01-somma-dispari}
Dimostrare che, per ogni \(n\in\N\) con \(n\ge1\), vale
\[
  \sum_{k=1}^{n}(2k-1)=n^2.
\]
\end{exercise}

\begin{proof}[Soluzione dell'esercizio]
Si applica il \cref{thm:principio-induzione} con valore iniziale \(p=1\).

\emph{Passo base.}
Per \(n=1\),
\[
  \sum_{k=1}^{1}(2k-1)=1=1^2.
\]

\emph{Passo induttivo.}
Sia \(n\ge1\) e si supponga che valga l'ipotesi induttiva
\[
  \sum_{k=1}^{n}(2k-1)=n^2.
\]
Separando l'ultimo addendo della somma fino a \(n+1\) e applicando
l'ipotesi induttiva, si ottiene
\begin{align*}
  \sum_{k=1}^{n+1}(2k-1)
    &=\sum_{k=1}^{n}(2k-1)+\bigl(2(n+1)-1\bigr)\\
    &=n^2+2n+1\\
    &=(n+1)^2.
\end{align*}
È dunque verificata la formula per \(n+1\). Per il principio di
induzione, l'identità vale per ogni \(n\ge1\).
\end{proof}

\begin{exercise}[Somma dei quadrati dei primi \(n\) naturali positivi]
\label{ex:tutorato01-somma-quadrati}
Dimostrare che, per ogni \(n\in\N\) con \(n\ge1\), vale
\[
  \sum_{k=1}^{n}k^2=\frac{n(n+1)(2n+1)}6.
\]
\end{exercise}

\begin{proof}[Soluzione dell'esercizio]
Si applica il \cref{thm:principio-induzione} con valore iniziale \(p=1\).

\emph{Passo base.}
Per \(n=1\),
\[
  \sum_{k=1}^{1}k^2=1=\frac{1\cdot2\cdot3}{6}.
\]

\emph{Passo induttivo.}
Sia \(n\ge1\) e si supponga
\[
  \sum_{k=1}^{n}k^2=\frac{n(n+1)(2n+1)}6.
\]
Separando l'ultimo addendo e applicando l'ipotesi induttiva,
\begin{align*}
  \sum_{k=1}^{n+1}k^2
    &=\sum_{k=1}^{n}k^2+(n+1)^2\\
    &=\frac{n(n+1)(2n+1)}6+(n+1)^2\\
    &=\frac{n+1}{6}\bigl(n(2n+1)+6(n+1)\bigr)\\
    &=\frac{n+1}{6}(2n^2+7n+6)\\
    &=\frac{(n+1)(n+2)(2n+3)}6.
\end{align*}
Nell'ultimo passaggio è stata utilizzata la fattorizzazione
\(2n^2+7n+6=(n+2)(2n+3)\). Poiché \(2n+3=2(n+1)+1\),
l'espressione ottenuta coincide con quella richiesta per \(n+1\).
Il passo induttivo è quindi verificato, e la formula vale per ogni
\(n\ge1\).
\end{proof}

\begin{exercise}[Dimostrazione del binomio di Newton]
\label{ex:tutorato01-binomio-newton}
Dimostrare per induzione il \cref{thm:binomio-newton}: per ogni
\(a,b\in\R\) e per ogni \(n\in\N\),
\[
  (a+b)^n=\sum_{k=0}^{n}\binom{n}{k}a^k b^{n-k}.
\]
\end{exercise}

\begin{proof}[Soluzione dell'esercizio]
Fissati \(a,b\in\R\), si applica il \cref{thm:principio-induzione}
con valore iniziale \(p=0\). Si mantiene la convenzione \(t^0=1\),
anche per \(t=0\), adottata nel \cref{thm:binomio-newton}.

\emph{Passo base.}
Per la \cref{def:coefficiente-binomiale}, \(\binom{0}{0}=1\).
Quindi, per \(n=0\),
\[
  (a+b)^0=1=\binom{0}{0}a^0b^0.
\]

\emph{Passo induttivo.}
Sia \(n\in\N\) e si supponga che valga l'ipotesi induttiva
\[
  (a+b)^n=\sum_{k=0}^{n}\binom{n}{k}a^k b^{n-k}.
\]
Moltiplicando per \(a+b\) e distribuendo il prodotto rispetto alla
somma finita, si ottiene
\begin{align*}
  (a+b)^{n+1}
    &=(a+b)^n(a+b)\\
    &=\left(\sum_{k=0}^{n}\binom{n}{k}a^k b^{n-k}\right)(a+b)\\
    &=\sum_{k=0}^{n}\binom{n}{k}a^{k+1}b^{n-k}
      +\sum_{k=0}^{n}\binom{n}{k}a^k b^{n+1-k}.
\end{align*}
Nella prima somma si pone \(j=k+1\). Quando \(k\) varia da \(0\)
a \(n\), il nuovo indice \(j\) varia da \(1\) a \(n+1\);
inoltre \(k=j-1\) e \(n-k=n+1-j\). Pertanto
\[
  \sum_{k=0}^{n}\binom{n}{k}a^{k+1}b^{n-k}
  =\sum_{j=1}^{n+1}\binom{n}{j-1}a^j b^{n+1-j}.
\]
Rinominando \(j\) come \(k\), l'uguaglianza precedente dà
\[
  (a+b)^{n+1}
  =\sum_{k=1}^{n+1}\binom{n}{k-1}a^k b^{n+1-k}
    +\sum_{k=0}^{n}\binom{n}{k}a^k b^{n+1-k}.
\]
Le due somme hanno ora gli stessi fattori \(a^k b^{n+1-k}\),
ma intervalli di sommazione diversi: gli indici comuni sono
\(1,\ldots,n\). Il termine con \(k=n+1\) compare soltanto nella
prima somma e vale \(a^{n+1}\), mentre quello con \(k=0\) compare
soltanto nella seconda e vale \(b^{n+1}\), poiché
\(\binom{n}{n}=\binom{n}{0}=1\).
Separando questi termini e combinando le somme sugli indici comuni,
si ottiene
\[
  (a+b)^{n+1}
  =b^{n+1}
    +\sum_{k=1}^{n}
      \left(\binom{n}{k-1}+\binom{n}{k}\right)a^k b^{n+1-k}
    +a^{n+1}.
\]
Per \(n=0\) non vi sono indici comuni: la somma da \(1\) a \(0\)
è vuota e vale \(0\), e l'uguaglianza resta valida.

Per ogni \(1\le k\le n\), l'identità di Pascal, dimostrata nella
\cref{prop:identita-pascal}, si applica con indice superiore \(n+1\)
e fornisce
\[
  \binom{n}{k-1}+\binom{n}{k}=\binom{n+1}{k}.
\]
Infatti, \(n+1\ge1\) e \(1\le k\le n+1\), come richiesto
dalle ipotesi dell'identità. Per la
\cref{def:coefficiente-binomiale}, anche i termini estremi hanno
coefficienti \(\binom{n+1}{0}=\binom{n+1}{n+1}=1\). Ne segue
\begin{align*}
  (a+b)^{n+1}
    &=b^{n+1}+\sum_{k=1}^{n}\binom{n+1}{k}a^k b^{n+1-k}
      +a^{n+1}\\
    &=\sum_{k=0}^{n+1}\binom{n+1}{k}a^k b^{n+1-k}.
\end{align*}
È quindi verificata la formula per \(n+1\). Il principio di induzione
assicura che essa vale per ogni \(n\in\N\). Poiché \(a,b\in\R\)
erano arbitrari, segue il \cref{thm:binomio-newton}.
\end{proof}
